\documentclass[11pt]{article}

\usepackage[margin=1in]{geometry}
\usepackage{graphicx}
\usepackage{booktabs}
\usepackage{caption}
\usepackage{subcaption}
\usepackage{amsmath}
\usepackage{amssymb}
\usepackage{hyperref}
\usepackage{xcolor}
\usepackage{placeins}
\usepackage{array}
\setlength{\emergencystretch}{2em}
\newcolumntype{L}[1]{>{\raggedright\arraybackslash}p{#1}}

\hypersetup{
  colorlinks=true,
  linkcolor=black,
  citecolor=black,
  urlcolor=blue
}

\graphicspath{{../outputs/figures/}}

\title{A Public-Data Bayesian Framework for External Wildfire Exposure Ranking of Transmission-Line Segments}

\author{Wenxing Yi}
\date{\today}

\begin{document}
\maketitle

\begin{abstract}
Linear energy infrastructure increasingly requires segment-level prioritization
under wildfire exposure, but many asset-condition and outage datasets are not
public. This paper presents a reproducible public-data framework for external
wildfire exposure ranking of California transmission-line segments. The lines
are treated as receptor assets rather than ignition sources. Using California
Energy Commission transmission lines, CAL FIRE FRAP fire perimeters and cause
fields, gridMET weather summaries, LANDFIRE LF2024 fuel and canopy layers, and
USGS 3DEP terrain, the workflow constructs a 2017--2023 segment-year panel for
the Northern Sierra and Southern Cascades. The main validation label excludes
CAL FIRE CAUSE = 11, Electrical Power, to reduce conflation between external
fire exposure and electrical-power fire cause. A Bayesian logistic exposure
model is compared with a deterministic public physics score and machine-learning
baselines under a 2022--2023 temporal holdout. The Bayesian and regularized
logistic models have similar discrimination, while the Bayesian model provides
posterior exceedance probabilities and posterior top-decile rank probabilities
for decision-facing triage. The framework is not an ignition, equipment-failure,
outage, or responsibility model.
\end{abstract}

\section{Introduction}

Wildfire exposure increasingly affects linear infrastructure systems. For
transmission corridors, planning and inspection decisions are often made at the
line-segment scale, before the exact location and severity of future fires are
known. Screening-level tools therefore need to prioritize segments using
weather, terrain, fuel, and historical fire information while keeping the target
definition clear.

Powerline-wildfire research can be separated into two related but distinct
directions. One direction studies grid-to-fire processes, such as
powerline-caused ignition. A second direction studies fire-to-grid exposure,
damage, and service disruption. This paper addresses only a limited part of the
second direction: external wildfire exposure ranking for transmission-line
segments treated as receptor assets. It does not estimate powerline-caused
ignition probability, equipment failure, outage probability, damage, or legal
responsibility.

The public-data constraint is central to the design. Private utility datasets
such as relay operations, outage logs, conductor type, asset age, inspection
records, vegetation-management histories, and equipment condition are not used.
The resulting target is therefore a screening-level external exposure ranking,
not a full operational risk model.

Existing deterministic hazard scores and machine-learning classifiers commonly
return a single score or probability. Rare-disaster triage also needs
uncertainty-aware quantities that answer decision questions directly, such as
whether a segment exceeds a probability threshold or whether it plausibly belongs
in the top decile of a ranked list. This motivates a Bayesian decision layer on
top of public fuel, weather, terrain, and historical exposure covariates.

The contributions are:
\begin{enumerate}
  \item a reproducible public-data segment-year exposure panel for California
  transmission-line segments;
  \item cause-screened external wildfire exposure labels that exclude
  electrical-power fire causes from the main validation target;
  \item public fuel, weather, and terrain covariate construction for
  segment-level ranking;
  \item Bayesian posterior exceedance and posterior top-decile rank
  probabilities as decision-facing outputs;
  \item temporal-holdout comparison with deterministic and machine-learning
  baselines; and
  \item an initial Stan MCMC check of the Laplace posterior prototype.
\end{enumerate}
\autoref{fig:advantage} summarizes this positioning relative to deterministic
scores and flexible ML baselines.

\begin{figure}[!htbp]
  \centering
  \includegraphics[width=0.95\textwidth]{phase8/phase8_research_advantage_matrix.png}
  \caption{Positioning of the proposed public-data exposure-ranking framework.}
  \label{fig:advantage}
\end{figure}

\section{Study Area and Public Data}

\subsection{Study Area}

The case study focuses on the Northern Sierra and Southern Cascades in
California. The decision units are approximately 1 km transmission-line
segments. The analysis panel spans 2017--2023 and uses segment-year rows so
that annual fire labels and annual fire-weather summaries are aligned. The
current prototype uses a fixed 1 km segment buffer for perimeter-overlap labels;
buffer sensitivity is left for follow-up validation.
\autoref{fig:study-area} shows the study region and receptor-asset segments.

\begin{figure}[!htbp]
  \centering
  \includegraphics[width=0.95\textwidth]{sci_phase1/fig1_study_area_transmission_fire_labels.png}
  \caption{Study region, transmission-line segments, and annual exposure labels.}
  \label{fig:study-area}
\end{figure}

\subsection{Public Data Sources}

\autoref{tab:data-sources} summarizes the public datasets and their roles. The
processed panel contains 11,255 line segments and 78,785 segment-year
observations. The environmental covariate layers used for quality control and
feature construction are shown in \autoref{fig:gridmet-layers},
\autoref{fig:dem-layers}, and \autoref{fig:landfire-layers}.

\begin{table}[!htbp]
  \centering
  \footnotesize
  \setlength{\tabcolsep}{2pt}
  \caption{Public data sources and roles.}
  \label{tab:data-sources}
  \begin{tabular}{L{0.13\textwidth} L{0.15\textwidth} L{0.18\textwidth} L{0.11\textwidth} L{0.16\textwidth} L{0.18\textwidth}}
    \toprule
    Dataset & Source & Spatial/temporal resolution & Years / version & Role & Notes and limitations \\
    \midrule
    Transmission lines & California Energy Commission & Line geometries & Current public release & Receptor assets & External exposure receptors; TODO citation. \\
    Fire perimeters & CAL FIRE FRAP & Perimeter polygons with cause fields & 2017--2023 subset & Annual exposure labels & Exposure proxy, not damage or outage; TODO citation. \\
    gridMET & University of Idaho gridMET & Daily, about 4 km; annual summaries & 2017--2023 & Weather covariates & Annual/fire-season summaries; TODO citation. \\
    LANDFIRE & LANDFIRE LF2024 & Fuel and canopy rasters & LF2024 & Fuel-structure covariates & Test time-matched releases; TODO citation. \\
    USGS 3DEP DEM & USGS 3DEP & Elevation and terrain rasters & Current public DEM & Terrain context & Not a fire-spread simulation; TODO citation. \\
    \bottomrule
  \end{tabular}
\end{table}

\begin{figure}[!htbp]
  \centering
  \includegraphics[width=0.95\textwidth]{data_previews/phase2_gridmet_preview_2021.png}
  \caption{gridMET annual fire-weather covariate layers for 2021.}
  \label{fig:gridmet-layers}
\end{figure}

\begin{figure}[!htbp]
  \centering
  \includegraphics[width=0.95\textwidth]{data_previews/phase2_dem_preview.png}
  \caption{USGS 3DEP terrain covariate layers.}
  \label{fig:dem-layers}
\end{figure}

\begin{figure}[!htbp]
  \centering
  \includegraphics[width=0.95\textwidth]{data_previews/phase2_landfire_preview.png}
  \caption{LANDFIRE LF2024 fuel and canopy covariate layers.}
  \label{fig:landfire-layers}
\end{figure}
\FloatBarrier

\section{Label Construction}

Let $i$ index line segments and $t$ index years. The main validation label is
\begin{equation}
y^{exo}_{it} =
\mathbb{1}\{\textrm{buffer}_i \cap \textrm{fire perimeter}_t \ne \emptyset,
\textrm{CAUSE} \ne 11\},
\end{equation}
where CAL FIRE CAUSE = 11 denotes Electrical Power. This exclusion is a label
boundary choice: it is used to avoid mixing external wildfire exposure with
electrical-power fire cause in the main validation target.

Four annual labels are retained. The all-fire label, $y^{all}_{it}$, records any
buffer-perimeter overlap. The exogenous label, $y^{exo}_{it}$, excludes CAUSE =
11 and is the main target. The strict label, $y^{strict}_{it}$, excludes CAUSE
in $\{11,2\}$, where CAUSE = 2 corresponds to Equipment Use. The
electrical-power overlap label, $y^{elec}_{it}$, records CAUSE = 11 overlap as
a diagnostic only. It is not interpreted as a calibrated powerline ignition
label because a perimeter-level cause code does not identify the exact ignition
asset, device, or line segment.

Across 2017--2023, the panel contains 3,187 all-fire exposure positives, 2,331
exogenous exposure positives, 2,075 strict exposure positives, and 885
electrical-power overlap positives. \autoref{fig:cause-screening} maps the
cause-screened labels, and \autoref{fig:annual-labels} reports annual label
counts.

\begin{figure}[!htbp]
  \centering
  \includegraphics[width=0.95\textwidth]{sci_phase1/fig2_cause_screened_label_maps.png}
  \caption{Cause-screened exposure labels for receptor-asset validation.}
  \label{fig:cause-screening}
\end{figure}

\begin{figure}[!htbp]
  \centering
  \includegraphics[width=0.90\textwidth]{sci_phase1/fig3_annual_segment_year_label_counts.png}
  \caption{Annual segment-year label counts.}
  \label{fig:annual-labels}
\end{figure}
\FloatBarrier

\section{Methods}

\subsection{Segment Construction and Exposure Buffers}

CEC transmission-line geometries are cleaned, clipped to the study region, and
split into approximately 1 km decision units. Each segment is assigned an annual
record for 2017--2023. The current prototype uses a fixed 1 km buffer around
each segment when constructing fire-perimeter overlap labels. Future work should
test 250 m, 500 m, 1000 m, and 2000 m buffers.

\subsection{Environmental Covariates}

The feature set combines dynamic gridMET dryness and wind summaries, static
LANDFIRE fuel-structure variables, USGS terrain covariates, and simple
physics-guided engineered indices. In this paper, ``physics-guided'' means
feature engineering informed by known fire-weather, fuel, and terrain controls;
it is not a physical fire-spread simulation. The selected feature core includes
VPD, ERC, BI, wind speed, 100-hour and 1000-hour fuel moisture, minimum relative
humidity, precipitation, elevation, canopy cover, canopy height, canopy base
height, canopy bulk density, dryness index, fuel-structure index, and a
dryness-by-fuel interaction. Covariate distribution diagnostics and spatial
quality-control maps are shown in \autoref{fig:covariate-distributions} and
\autoref{fig:spatial-qc}.

\begin{figure}[!htbp]
  \centering
  \includegraphics[width=0.95\textwidth]{phase3/phase3_covariate_qc_distributions.png}
  \caption{Covariate distribution diagnostics.}
  \label{fig:covariate-distributions}
\end{figure}

\begin{figure}[!htbp]
  \centering
  \includegraphics[width=0.95\textwidth]{phase3/phase3_spatial_qc_maps.png}
  \caption{Spatial covariate quality-control maps.}
  \label{fig:spatial-qc}
\end{figure}

\subsection{Deterministic Comparator}

The deterministic comparator is a calibrated public physics-guided score based
on dryness, fuel structure, and elevation. It is used as a transparent benchmark
for ranking, not as the final decision model. Let $z_{k,it}$ denote standardized
covariates estimated using the training period. The deterministic baseline can
be written as
\begin{equation}
s_{it}^{det}
= \mathrm{clip}_{[0,1]}
\left[
a_1 z^{dry}_{it}
+ a_2 z^{fuel}_{i}
+ a_3 z^{elev}_{i}
\right].
\end{equation}
The score is subsequently calibrated to a probability scale using regularized
logistic calibration. It yields a single ranking but no posterior distribution
over ranks.

\subsection{Bayesian Exposure Model}

The main Bayesian model is a weakly regularized logistic exposure model. For
segment $i$ and year $t$:
\begin{align}
y^{exo}_{it} &\sim \mathrm{Bernoulli}(p_{it}), \\
\eta_{it} &= \alpha + \mathbf{x}_{it}^{\top}\boldsymbol{\beta}, \\
p_{it} &= \mathrm{logit}^{-1}(\eta_{it}), \\
\alpha &\sim \mathcal{N}(0, 5^2), \\
\beta_k &\sim \mathcal{N}(0, 2^2).
\end{align}
The posterior is
\begin{equation}
p(\alpha,\boldsymbol{\beta}\mid D)
\propto
\prod_{(i,t)\in D}
p_{it}^{y_{it}}(1-p_{it})^{1-y_{it}}
p(\alpha)\prod_k p(\beta_k).
\end{equation}
The prototype is fit using a Laplace posterior approximation around the
posterior mode and checked with an initial Stan MCMC run.

Posterior draws are indexed by $s=1,\ldots,S$:
\begin{align}
(\alpha^{(s)},\boldsymbol{\beta}^{(s)})
&\sim p(\alpha,\boldsymbol{\beta}\mid D), \\
p_{it}^{(s)}
&= \mathrm{logit}^{-1}
\left(\alpha^{(s)}+\mathbf{x}_{it}^{\top}\boldsymbol{\beta}^{(s)}\right).
\end{align}
For segment-level decisions over the holdout period $T^\star$, posterior segment
exposure is aggregated as
\begin{equation}
\bar{p}_{i}^{(s)}
= \frac{1}{|T^\star|}
\sum_{t\in T^\star} p_{it}^{(s)}.
\end{equation}

The decision quantities are:
\begin{align}
q_i
&= P(\bar{p}_i > \tau \mid D)
\approx \frac{1}{S}\sum_{s=1}^{S}
\mathbb{1}\{\bar{p}_{i}^{(s)}>\tau\}, \\
\rho_i
&= P(\mathrm{rank}(\bar{p}_i) \le K \mid D)
\approx \frac{1}{S}\sum_{s=1}^{S}
\mathbb{1}\{\mathrm{rank}(\bar{p}_{i}^{(s)})\le K\},
\end{align}
where $\tau$ is the training prevalence and $K$ is the top-decile segment count.
\autoref{fig:method-schematic} summarizes this workflow.

\begin{figure}[!htbp]
  \centering
  \includegraphics[width=0.95\textwidth]{phase8/phase8_bayesian_method_schematic-2.png}
  \caption{Bayesian workflow for posterior exposure and rank probabilities.}
  \label{fig:method-schematic}
\end{figure}

\subsection{Validation Design}

Models are trained on 2017--2021 and evaluated on 2022--2023. The holdout set
contains 22,510 segment-year rows and 356 exogenous exposure positives, giving a
holdout prevalence of 1.58\%. This rare-event baseline is important for
interpreting PR-AUC. Metrics include ROC-AUC, PR-AUC, Brier score, row-level
top-decile capture, and segment-level top-decile capture.
\FloatBarrier

\section{Results}

\subsection{Temporal Holdout Performance}

\autoref{tab:main-comparison} compares the deterministic score, Bayesian
Laplace model, Stan MCMC pilot, and machine-learning baselines. The Bayesian and
regularized logistic L2 models have nearly identical ROC-AUC and PR-AUC. The
Bayesian contribution is therefore not a claim of broad classification
superiority; it is the posterior decision layer shown in
\autoref{fig:det-vs-bayes}, \autoref{fig:bayes-maps}, and
\autoref{fig:posterior-maps}. Relative to the deterministic calibrated score,
the Bayesian ranking improves rare-event decision concentration in this temporal
holdout, summarized in \autoref{fig:temporal-holdout}.

\begin{figure}[!htbp]
  \centering
  \includegraphics[width=0.95\textwidth]{phase4/phase4_temporal_holdout_validation.png}
  \caption{Temporal-holdout validation for deterministic and Bayesian ranking.}
  \label{fig:temporal-holdout}
\end{figure}

\begin{figure}[!htbp]
  \centering
  \includegraphics[width=0.95\textwidth]{phase4/phase4_deterministic_vs_bayesian_ranking.png}
  \caption{Deterministic score and Bayesian posterior decision metrics.}
  \label{fig:det-vs-bayes}
\end{figure}

\begin{figure}[!htbp]
  \centering
  \includegraphics[width=0.95\textwidth]{phase4/phase4_decision_metric_maps.png}
  \caption{Posterior decision metrics mapped to transmission-line segments.}
  \label{fig:bayes-maps}
\end{figure}

\begin{table}[!htbp]
  \centering
  \small
  \caption{Temporal-holdout comparison for 2022--2023 exogenous exposure.}
  \label{tab:main-comparison}
  \begin{tabular}{lrrrrr}
    \toprule
    Model & ROC-AUC & PR-AUC & Brier & Row top 10\% & Segment top 10\% \\
    \midrule
    ML logistic L2 & 0.663 & 0.039 & 0.015713 & 0.289 & 0.216 \\
    Bayesian Laplace & 0.662 & 0.038 & 0.015718 & 0.289 & 0.216 \\
    Stan MCMC pilot & 0.662 & 0.038 & 0.015710 & 0.289 & 0.216 \\
    Deterministic calibrated & 0.651 & 0.028 & 0.015734 & 0.219 & 0.118 \\
    Balanced extra trees & 0.617 & 0.022 & 0.056091 & 0.180 & 0.104 \\
    Histogram gradient boosting & 0.553 & 0.019 & 0.016647 & 0.132 & 0.149 \\
    Balanced random forest & 0.568 & 0.018 & 0.034539 & 0.107 & 0.152 \\
    \bottomrule
  \end{tabular}
\end{table}

\subsection{Feature Ablation and Baseline Checks}

\autoref{tab:ablation} and \autoref{fig:ablation} summarize feature-set
ablations. Dynamic gridMET weather summaries carry much of the temporal-holdout
discrimination. LANDFIRE fuel variables are useful context variables but do not
by themselves support strong holdout discrimination in this prototype.
\autoref{fig:ml-comparison} compares deterministic, Bayesian, and ML baselines,
and \autoref{fig:calibration} reports the calibration diagnostic.

\begin{figure}[!htbp]
  \centering
  \includegraphics[width=0.95\textwidth]{phase5/phase5_ablation_validation_metrics.png}
  \caption{Feature-set ablation under temporal holdout.}
  \label{fig:ablation}
\end{figure}

\begin{table}[!htbp]
  \centering
  \small
  \caption{Feature-set ablation under 2022--2023 temporal holdout.}
  \label{tab:ablation}
  \begin{tabular}{lrrrrr}
    \toprule
    Model & ROC-AUC & PR-AUC & Brier & Row top 10\% & Segment top 10\% \\
    \midrule
    M0 geometry + terrain & 0.489 & 0.024 & 0.016052 & 0.157 & 0.157 \\
    M1 + LANDFIRE fuel & 0.543 & 0.020 & 0.016076 & 0.149 & 0.149 \\
    M2 + gridMET weather & 0.665 & 0.029 & 0.015822 & 0.199 & 0.143 \\
    M3 fuel + weather & 0.667 & 0.030 & 0.015821 & 0.222 & 0.152 \\
    M4 physics interaction & 0.671 & 0.030 & 0.015820 & 0.211 & 0.152 \\
    M5 selected physics core & 0.662 & 0.038 & 0.015723 & 0.289 & 0.219 \\
    \bottomrule
  \end{tabular}
\end{table}

\begin{figure}[!htbp]
  \centering
  \includegraphics[width=0.95\textwidth]{phase6/phase6_model_comparison_metrics.png}
  \caption{Deterministic, Bayesian, and machine-learning baseline comparison.}
  \label{fig:ml-comparison}
\end{figure}

\begin{figure}[!htbp]
  \centering
  \includegraphics[width=0.85\textwidth]{phase6/phase6_calibration_check.png}
  \caption{Calibration diagnostic for rare external exposure probabilities.}
  \label{fig:calibration}
\end{figure}

\begin{figure}[!htbp]
  \centering
  \includegraphics[width=0.95\textwidth]{phase8/phase8_faceted_zoom_posterior_maps.png}
  \caption{Full-region and zoomed maps of posterior decision metrics.}
  \label{fig:posterior-maps}
\end{figure}
\FloatBarrier

\section{Discussion}

The main result is a bounded one: public data can support a reproducible,
screening-level framework for external wildfire exposure ranking of
transmission-line receptor assets. The output should be read as a ranked
exposure screen, not as ignition causality, equipment condition, outage
likelihood, damage, or responsibility.

The temporal-holdout results do not show that Bayesian modeling is universally
more accurate than all machine-learning alternatives. The regularized logistic
L2 baseline and the Bayesian Laplace model have nearly identical discrimination
in \autoref{tab:main-comparison}. The Bayesian advantage in this paper is
decision-facing: it converts the fitted exposure model into posterior
exceedance probability and posterior top-decile rank probability. These
quantities are more aligned with rare-event triage than a single deterministic
score because they represent uncertainty in threshold and ranking decisions.

The deterministic comparator remains useful because it is transparent and easy
to audit. Its limitation is that it produces one calibrated score and therefore
cannot express posterior rank uncertainty. Tree-based ML baselines are more
flexible, but in this temporal split they do not improve rare-event
concentration relative to the selected physics-guided linear core. This result
should be interpreted as evidence for cautious feature design and temporal
validation, not as a general rejection of ML models.

Existing wildfire risk tools often target different questions. Powerline
ignition and equipment-failure models generally require private utility data and
are aimed at grid-to-fire causality. Fire-spread simulators can represent richer
physical dynamics but require scenario assumptions and are not automatically
calibrated for public-data segment ranking. The present framework occupies a
narrower space: fire-to-grid external exposure ranking using public data and a
posterior decision layer.

\section{Limitations}

Several limitations are important for an EarthArXiv preprint. First, the
pipeline estimates external exposure only. It does not estimate ignition
causality, equipment failure, outage probability, damage, or legal
responsibility. Second, the CAL FIRE cause field is a perimeter-level attribute;
electrical-power overlap is therefore retained only as a diagnostic label.
Third, the current prototype uses LANDFIRE LF2024 static fuel layers for a
2017--2023 panel. This may create temporal mismatch with earlier years and
should be tested against time-matched LANDFIRE releases in a journal version.
Fourth, the label construction uses a fixed 1 km buffer. Buffer sensitivity
should compare 250 m, 500 m, 1000 m, and 2000 m definitions. Fifth, temporal
holdout is not a substitute for spatial block validation. Finally, the Stan run
is an initial MCMC check; full Stan diagnostics, longer chains, convergence
summaries, and posterior predictive checks are still needed.

\section{Conclusion}

This paper presents a public-data framework for external wildfire exposure
ranking of transmission-line segments treated as receptor assets. The workflow
builds a 2017--2023 segment-year panel, uses cause-screened exposure labels, and
compares deterministic, Bayesian, and machine-learning baselines under a
2022--2023 temporal holdout. The Bayesian and regularized logistic models show
similar discrimination, while the Bayesian formulation adds posterior
exceedance and posterior top-decile rank probabilities. These posterior decision
quantities are the main methodological contribution for rare-event exposure
triage. The framework should not be interpreted as a powerline ignition,
equipment-failure, outage, damage, or responsibility model.

\appendix
\section{Reproducibility Notes}

All data processing, feature construction, model fitting, validation, and figure
generation are scripted in the project repository. The public-data workflow is
organized by phases: asset and label construction, environmental covariate
preparation, model-table construction, Bayesian ranking, ablation, ML/Stan
validation, and manuscript result generation. Manuscript tables and figures are
generated from scripted outputs rather than manually transcribed. Public-source
citations remain marked as TODO placeholders until a bibliography is added.

\end{document}
